% Venue-neutral body for the signal-liveness workshop paper.
% No documentclass, no author block, no bibliography command: those
% live in the per-venue wrappers, exactly as in the survey.
%
% Sections 1, 2, 3 and 7 are drafted. Sections 4, 5 and 6 are scaffolded
% and cannot be written until the measurements exist. See
% SKELETON-REVIEW.md for why the existing pilot run does not serve.

\section{Introduction}
\label{sec:intro}

An LLM inference cluster under overload has one shared postmortem shape.
Utilization climbs, a queue forms, latency degrades before errors appear, and
the operator discovers that the system accepted more work than it could finish.
Recent public incidents name the same layer as the failure site: requests were
admitted that should have been refused, and the refusal, when it came, came from
a timeout rather than a decision.

Admission control is the decision taken before a request enters the system:
admit it, hold it, or refuse it. A companion systematic review examined
fifty-five admission and admission-adjacent works published between 2022 and
2026 and classified them on five axes \citep{\surveyref}. Applying the taxonomy
to the sixteen admission-primary works in that corpus, the review recorded seven
properties per paper: the five axes, plus what the policy assumes an arriving
request carries, plus the range the policy's signal took during the paper's own
evaluation. Six of the seven fill from published text, and all five axes are
among them.

One property fills for no paper. For none of the sixteen could the reviewers
determine the range the policy's own signal took during its evaluation. This
matters because of what it makes indistinguishable. A threshold that was never
approached and a threshold that was approached but not crossed produce the same
admission counts and the same downstream latencies. No reported number separates
a policy that was exercised from one that was merely present, which means a
reported improvement cannot be attributed to the policy's decisions rather than
to its inactivity.

The review states plainly that this finding is bibliographic rather than
empirical: it reports an absence in what was extracted, not an observation of a
policy failing this way in a deployment. Establishing the latter requires
running a policy whose signal is known to move alongside one whose signal is
known not to, on the same deployment and against the same trace, and asking
whether any reported number separates them. The review states this as an open
problem.

This paper is an attempt at that open problem. We contribute:

\begin{enumerate}
  \item A \textbf{signal-liveness probe}: a protocol that takes a policy's
    declared signal and a target deployment and reports whether the signal
    reaches its decision threshold under offered load, runnable in an afternoon
    on a single GPU.
  \item \textbf{Measurements} of two published policies on unmodified vLLM
    across arrival regimes, reporting signal range alongside the admission
    decisions it produced, and quantifying how far two policies can diverge on
    individual requests while agreeing on the aggregate rate that papers
    conventionally report.
  \item A \textbf{capability-declaration schema} derived from the taxonomy,
    letting a policy state the signal it reads and the band it assumes, and a
    runtime gate that refuses deployment when the engine cannot supply the
    signal or when prior probes show the signal pinned in the target regime.
\end{enumerate}

This is a methodological contribution rather than a new admission mechanism. We
propose no policy and claim no improvement over any published one. The precedent
for evaluation-harness work at a systems venue is MLPerf Inference
\citep{reddi2020mlperf}.

\section{Background: the taxonomy and the unfilled property}
\label{sec:background}

The taxonomy classifies an admission policy on five axes, each answerable about
one's own deployment without reading the papers: the \emph{unit} at which
admission is decided, the \emph{setting} the policy assumes, the
\emph{objective} it optimizes, the \emph{quantity} it watches, and its
\emph{portability}, meaning whether it can run in front of an unmodified engine
or requires control over batch formation.

All five are recoverable from published text for all sixteen admission-primary
works. Two further properties were recorded and are not axes. The first, what a
policy assumes an arriving request carries, fills for every paper. The second,
the range the signal took during evaluation, fills for none.

\begin{table}[t]
\centering
\caption{The seven properties recorded per paper in the companion review, and
whether each is recoverable from published text across the sixteen
admission-primary works.}
\label{tab:properties}
\small
\begin{tabular}{llc}
\toprule
\textbf{Property} & \textbf{Kind} & \textbf{Recoverable} \\
\midrule
Unit of admission        & axis     & 16/16 \\
Setting                  & axis     & 16/16 \\
Objective                & axis     & 16/16 \\
Signal quantity          & axis     & 16/16 \\
Portability              & axis     & 16/16 \\
Request metadata assumed & recorded & 16/16 \\
\textbf{Signal range observed} & \textbf{recorded} & \textbf{0/16} \\
\bottomrule
\end{tabular}
\end{table}

We adopt the following definition, which the companion review does not state
formally and which is this paper's own.

\begin{definition}[Signal liveness]
\label{def:liveness}
Let $s(t)$ be the quantity a policy reads and $\theta$ its decision threshold.
Given an offered-load probe over interval $[0,T]$ and a band width
$\delta$, the signal is \textbf{live} if the fraction of samples with
$|s(t) - \theta| < \delta$ is non-negligible, and \textbf{inert} otherwise.
A policy whose signal is inert cannot have acted on its threshold during the
run, whatever its reported admission rate.
\end{definition}

The definition is deliberately relative to a deployment rather than to a policy.
Whether a signal is live depends on the engine configuration, the model, and the
offered load, and one of this paper's results is that the same policy reading
the same quantity can be live in one configuration and inert in another.

\section{Method: the signal-liveness probe}
\label{sec:method}

The probe has three steps and requires no engine modification.

\paragraph{Declare capability.} The policy states the signal it reads, drawn
from the quantities the taxonomy enumerates (KV pressure, queue depth, deadline
slack, tenant credit, session state, wait estimate), together with its decision
threshold $\theta$ and the band $\delta$ within which it considers the signal
actionable.

\paragraph{Offer load.} The deployment is driven at a set of arrival rates
spanning its rated capacity, established by a calibration run rather than
assumed. Establishing capacity first is not optional: a probe run below the
saturation point reports a flat signal for every policy, and reads as inertness
when it is only absence of load.

\paragraph{Record range.} The declared signal is sampled throughout at fixed
cadence, and we report its minimum, maximum, median and upper percentiles,
together with the fraction of samples inside the policy's band and the admission
decisions taken at each sample.

The reported artifact is the pair: signal range and admission outcome. Either
alone is what the literature already provides and what
Definition~\ref{def:liveness} shows to be insufficient.

\section{Experimental setup}
\label{sec:setup}

% PENDING. Do not write against the existing pilot configuration: a 0.5B model
% given ~19 GB of KV cache with max_num_seqs=10 cannot fill that cache, so KV
% utilization near zero is arithmetic rather than a finding. The setup must
% first establish a configuration in which the signal can reach its threshold.
% See SKELETON-REVIEW.md.

\section{Results}
\label{sec:results}

% PENDING measurements. Structure: signal range per (policy, regime) cell as the
% central table; per-decision agreement between policies as the figure that
% carries the argument; the crossover point at which a pinned signal becomes
% live, if one exists in the tested range.

\section{Discussion}
\label{sec:discussion}

% PENDING results. Three parts: what one added column to a results table would
% have shown; the capability-declaration schema and runtime gate; and an
% explicit statement of what two policies on one engine cannot establish about
% the literature as a whole.

\section{Threats to validity}
\label{sec:threats}

\paragraph{Two policies is not a survey of policies.} We measure two of the
sixteen admission-primary policies in the corpus. This cannot establish how
often inert signals occur in the literature. It can establish that inertness is
measurable, that it occurs for at least one published policy in a configuration
a practitioner might plausibly choose, and that its absence from reported
results is a genuine information gap rather than a stylistic omission.

\paragraph{One engine.} We measure on vLLM. Other engines expose different
signals at different cadences, and a signal that is live on one may be
unobservable on another. The probe protocol is engine-agnostic; the measurements
are not.

\paragraph{Liveness is configuration-dependent, including ours.} The central
claim of this paper is that signal liveness depends on the deployment rather
than on the policy. That claim applies reflexively. Our own configuration is one
point in a space, and a reader should treat our numbers as an existence proof
about that point rather than as a characterization of the policy.

\paragraph{Synthetic workloads.} Our arrival processes and request shapes are
proxies for production traffic. Agentic and multi-tenant workloads have
correlation structure that a Poisson or bursty synthetic does not reproduce, and
that structure is exactly what drives the signals under study.

\section{Conclusion}
\label{sec:conclusion}

% PENDING results. Close on the reporting ask, which is one column, and on
% one-command reproduction from a public commit.
